\section{Locally compact modules and the Artin map}\label{sec:lca}
In this section, we compare the $K$-theory of locally compact $D$-modules with that of locally compact $F$-modules and relate this comparison to the Artin map.
The motivation comes from Clausen's description of the Artin map of class field theory through the $K$-theory of locally compact $F$-modules \cite{ClausenArtin}, and from its noncommutative analogue by Braunling--Henrard--van Roosmalen \cite{BHRAdeles}, which describes the id\`ele class group of $D$ modulo the id\`eles of local reduced norm one through $K_1$ of locally compact $D$-modules.
Braunling, Henrard and van Roosmalen conjecture that $K_2(\operatorname{LCA}_A)\cong K_2(\operatorname{LCA}_F)$ for every finite-dimensional simple $\Q$-algebra $A$ with center $F$ \cite[Conjecture 10.2]{BHRAdeles}.
For $A=D$, Theorem~\ref{thm:paper-lca-k2} below gives this isomorphism with integral coefficients, also when $D$ is ramified at a real place.

Our comparison shows that, in positive degrees, the fiber of the comparison $c^{\mathrm{lc}}_D$ is exactly the suspension $\Sigma Z$ of the global comparison fiber of Section~\ref{sec:global-spectrum}, a finite sum of shifted Eilenberg--Mac Lane spectra over number fields, with $2$ inverted under real ramification, and that this fiber is invisible to the Artin map.
On $\pi_1$ the compatibility gives the formula for the pro-$p$ Artin map on id\`ele classes: local reduced norms, followed by the global Artin map and the projection to the maximal pro-$p$ quotient, with $p$ odd under real ramification.

The $K$-theory of locally compact modules fits into a localization sequence between $K(D)$ and the restricted product of the local $K$-theories, and a comparison square built from a common integral model of $D$ induces a comparison $c^{\mathrm{lc}}_D$ between the cofibers.
If the global comparison of Section~\ref{sec:global-spectrum} is kept fixed, this square may have to be corrected.

Throughout this section, $F$ is a number field and, by Morita equivalence, $D$ is a central division algebra over $F$.
We work with integral coefficients when $D$ splits at every real place.
If $D$ is ramified at a real place, all spectra and groups in this section, except in Theorem~\ref{thm:paper-lca-k2}, are taken with $\Z[1/2]$ coefficients; the restricted products below are then formed integrally and $2$ is inverted afterwards, not place by place.
We construct $c^{\mathrm{lc}}_D$ (Theorem~\ref{thm:paper-lca-comparison}), identify its fiber with the suspension of the global comparison fiber (Theorem~\ref{thm:paper-lca-fiber-splitting}), prove the integral isomorphism on $K_2$ (Theorem~\ref{thm:paper-lca-k2}), and show that it is compatible with Clausen's Artin map (Theorem~\ref{thm:paper-lca-artin}).

\subsection{Locally compact modules and localization in positive degrees}
\label{subsec:paper-lca-localization}
We first set up the localization fiber sequences for locally compact modules over $D$ and over $F$ in positive degrees.
The input is the theorem of Braunling--Henrard--van Roosmalen, which we also need for second-countable modules.

For $R=D,F$, let $\operatorname{LCA}_R$ denote the category of locally compact right modules over the discrete ring $R$, with continuous $R$-linear maps.
We give it the standard exact structure in which closed embeddings are admissible monomorphisms and open surjections are admissible epimorphisms.
Write $\operatorname{AB}_R=\operatorname{LCA}_{R,\mathrm{ab}}$ for its full exact subcategory of adelic blocks.
Concretely, $\operatorname{AB}_R$ is the filtered $2$-colimit of the following categories.
For a maximal order $\mathfrak R$ in $R$ and a finite set $S$ of finite places, put
\[
 \mathcal J_R^S=
 \prod_{v\mid\infty}\operatorname{proj}(R_v)
 \times\prod_{v\in S}\operatorname{proj}(R_v)
 \times\prod_{v\notin S,\ v\nmid\infty}
                       \operatorname{proj}(\mathfrak R_v).
\]
For $S\subset S'$, the transition functor is extension of scalars at the newly added places.
This description is independent of the maximal order.
We use the theorem of Braunling--Henrard--van Roosmalen in the following form \cite[Definition 4.3, Proposition 4.5, Theorems 6.1, 6.2 and 6.13]{BHRAdeles}.

\begin{theorem}\label{thm:paper-lca-source}
There is a natural fiber sequence
\begin{equation}\label{eq:paper-lca-bhr}
 K(R)\xrightarrow{\lambda_R}K(\operatorname{AB}_R)
       \longrightarrow K(\operatorname{LCA}_R).
\end{equation}
The first map is induced by $P\mapsto P\otimes_R\mathbb A_R$.
For the middle term, there is an equivalence of spectra
\begin{equation}\label{eq:paper-lca-restricted-spectrum}
 K(\operatorname{AB}_R)\simeq
 \varinjlim_S\left(
 \prod_{v\mid\infty}K(R_v)\times\prod_{v\in S}K(R_v)
 \times\prod_{v\notin S,\ v\nmid\infty}K(\mathfrak R_v)\right).
\end{equation}
In particular, $K_j(\operatorname{AB}_R)=\prod'_vK_j(R_v)$, where the restriction is that the $v$-component lies in the image of $K_j(\mathfrak R_v)$ at all but finitely many finite places $v$.
\end{theorem}

The restricted product in~\eqref{eq:paper-lca-restricted-spectrum} is formed in spectra, using the product theorem for nonconnective $K$-theory of exact categories and the compatibility of $K$-theory with filtered colimits.

\begin{lemma}\label{lem:paper-lca-countable-localization}
Let $R\in\{D,F\}$, let $\mathcal C=\operatorname{LCA}_{R,\aleph}$ be the full exact subcategory of second-countable objects of $\operatorname{LCA}_R$, and let $\mathcal C_{\mathrm c}$ be its compact subcategory.
\begin{enumerate}
\item Every adelic block lies in $\mathcal C$.
\item The localization sequences of \cite[Proposition 6.3 and Corollary 6.6]{BHRAdeles}, for the quotient by $\mathcal C_{\mathrm c}$ and for the image $\mathcal V\simeq\operatorname{AB}_R$ of the adelic blocks in the exact hull $(\mathcal C/\mathcal C_{\mathrm c})^{\mathrm{ex}}$, hold with $\operatorname{LCA}_R$ replaced by $\mathcal C$.
\item The discrete quotient functor $M\mapsto M/Q(M)$ of \cite[Proposition 2.11]{BHRAdeles} induces an equivalence
\[
 (\mathcal C/\mathcal C_{\mathrm c})^{\mathrm{ex}}/\mathcal V
 \simeq\operatorname{Mod}_\aleph(R)/\operatorname{mod}(R),
\]
whose inverse is induced by the inclusion of countable discrete modules; here $\operatorname{Mod}_\aleph(R)$ is the category of countable discrete modules and $\operatorname{mod}(R)$ is the category of finitely generated modules.
\end{enumerate}
\end{lemma}
\begin{proof}
Second countability of locally compact groups is preserved under passage to closed subgroups, Hausdorff quotients, and finite iterated extensions.
A discrete quotient of a second-countable group is countable.
Since $R$ itself is countable, every finitely generated $R$-module is countable.
Every adelic block is second countable, since it is a restricted product of finite-dimensional local components over countably many places, so restricting to $\mathcal C$ does not change the category of adelic blocks.

We now explain how the localization proof of Braunling--Henrard--van Roosmalen restricts to this subcategory.
Let $M_0$ be an object of $\operatorname{LCA}_R$, and choose an order $\Lambda\subset R$.
As in the proof of \cite[Proposition~2.11]{BHRAdeles}, choose a conflation in $\operatorname{LCA}_\Lambda$
\[
 V\oplus C\hookrightarrow M_0\twoheadrightarrow M_0/(V\oplus C),
\]
where $V$ is a vector $R$-module, $C$ is a compact $\Lambda$-module that is open in a closed complement of $V$, and $M_0/(V\oplus C)$ is discrete, so that $V\oplus C$ is open in $M_0$; the summand $C$ need not be an $R$-module, but its saturation below is.
The split torsion pair of \cite[Proposition 2.11]{BHRAdeles} is constructed from the open $R$-submodule $Q(M_0)=\bigcup_{n\ge1}n^{-1}(V\oplus C)$ and the discrete quotient $M_0/Q(M_0)$.
If $M_0\in\mathcal C$, then $Q(M_0)$ is second countable and $M_0/Q(M_0)$ is countable.
The decomposition of \cite[Lemma 2.12]{BHRAdeles} splits the quasi-adelic part into a compact part and an adelic part by using only closed submodules and quotients, so both parts stay in $\mathcal C$.
Compact $R$-modules and adelic blocks are injective in the whole category and remain injective after restriction.

Consequently, we can apply to this subcategory the localization by the compact subcategory in \cite[Proposition 6.3]{BHRAdeles} and the full faithfulness of adelic blocks in \cite[Lemma 6.4]{BHRAdeles}.
Indeed, the quotient $\mathcal C/\mathcal C_{\mathrm c}$ and its exact hull are built intrinsically.
The category $\mathcal C$ has closed kernels, Hausdorff cokernels and finite pushouts and pullbacks, $\mathcal C_{\mathrm c}$ is a Serre subcategory, and images of maps from compact objects are closed and compact, so the first localization of \cite[\S5]{BHRAdeles} applies within $\mathcal C$; the arguments of \cite[Lemma 6.4, Proposition 6.5]{BHRAdeles} are then carried out with roofs in this quotient.
Let $\mathcal V$ denote the image of the adelic blocks in the exact hull $(\mathcal C/\mathcal C_{\mathrm c})^{\mathrm{ex}}$; then $\mathcal V\simeq\operatorname{AB}_R$.
The verification of percolation in \cite[Proposition 6.5]{BHRAdeles} is unchanged.
The main conditions of percolation use the fact that an admissible image of an adelic block is quasi-adelic and becomes an adelic block after removing its compact part, together with the injectivity of adelic blocks.
The remaining conditions involve only finitely many kernels, cokernels, and pushouts.
We therefore obtain the localization sequence of \cite[Corollary 6.6]{BHRAdeles} in this subcategory.

For~(3), we follow \cite[Proposition 6.8 and Corollary 6.9]{BHRAdeles}.
By \cite[Lemma 2.18]{BHRAdeles}, the discrete quotient functor $M\mapsto M/Q(M)$ is exact up to a finitely generated defect, so it induces an exact functor to $\operatorname{Mod}_\aleph(R)/\operatorname{mod}(R)$.
In the other direction, the inclusion of countable discrete modules kills $\operatorname{mod}(R)$ in $(\mathcal C/\mathcal C_{\mathrm c})^{\mathrm{ex}}/\mathcal V$, because for finitely generated projective $P$ the cokernel of $P\to P\otimes_R\mathbb A_R$ is compact.
The two functors are mutually inverse: a countable discrete module is its own discrete quotient, and $M\to M/Q(M)$ becomes an isomorphism once the compact and adelic parts of $Q(M)$ are removed.
None of these operations requires uncountable objects.
\end{proof}

\begin{lemma}\label{lem:paper-lca-countability}
Let $R\in\{D,F\}$.
The inclusion of the subcategory $\operatorname{LCA}_{R,\aleph}$ of Lemma~\ref{lem:paper-lca-countable-localization} induces an equivalence
\[
 K(\operatorname{LCA}_{R,\aleph})\xrightarrow{\ \sim\ }
 K(\operatorname{LCA}_R).
\]
Theorem~\ref{thm:paper-lca-source} also holds with $\operatorname{LCA}_R$ replaced by $\operatorname{LCA}_{R,\aleph}$.
\end{lemma}
\begin{proof}
Put $\mathcal C=\operatorname{LCA}_{R,\aleph}$, and write $\mathcal C_{\mathrm c}$ for its compact subcategory.
By Lemma~\ref{lem:paper-lca-countable-localization}, the localization of Braunling--Henrard--van Roosmalen and its discrete quotient restrict to $\mathcal C$.

Since $\mathcal C_{\mathrm c}$ is closed under countable products and $\operatorname{Mod}_\aleph(R)$ under countable sums, the Eilenberg swindle gives
\[
 K(\mathcal C_{\mathrm c})=0,
 \qquad K(\operatorname{Mod}_\aleph(R))=0.
\]
From the two quotient identifications and the localization sequences of Lemma~\ref{lem:paper-lca-countable-localization}, we obtain the same bicartesian square as in the proof of \cite[Theorem 6.2]{BHRAdeles}:
\[
\begin{array}{ccc}
 K(\operatorname{mod}(R))&\longrightarrow&K(\operatorname{AB}_R)\\
 \downarrow&&\downarrow\\
 K(\operatorname{Mod}_\aleph(R))&\longrightarrow&K(\mathcal C)
\end{array}
\]
The top arrow is induced by $P\mapsto P\otimes_R\mathbb A_R$.
The homotopy making the square commute comes from the exact sequence with a diagonal discrete lattice and a compact quotient
\[
 P\longrightarrow P\otimes_R\mathbb A_R
 \longrightarrow (P\otimes_R\mathbb A_R)/P.
\]
Since $K(\operatorname{Mod}_\aleph(R))=0$, the sequence $K(R)\to K(\operatorname{AB}_R)\to K(\mathcal C)$ is a fiber sequence.
The inclusion $\mathcal C\subset\operatorname{LCA}_R$ is compatible with the exact sequence above and induces a map of fiber sequences
\[
\begin{array}{ccccc}
 K(R)&\longrightarrow&K(\operatorname{AB}_R)&\longrightarrow&K(\mathcal C)\\
 \|&&\|&&\downarrow\\
 K(R)&\longrightarrow&K(\operatorname{AB}_R)&\longrightarrow&K(\operatorname{LCA}_R)
\end{array}
\]
whose bottom row is~\eqref{eq:paper-lca-bhr}.
Therefore the right vertical map is also an equivalence.
\end{proof}

Set
\[
 X_R=\tau_{\ge1}K(R),\qquad
 B_R=\tau_{\ge1}K(\operatorname{AB}_R),\qquad
 Y_R=\tau_{\ge1}K(\operatorname{LCA}_R).
\]
The map $K_0(R)\to K_0(\operatorname{AB}_R)$ is injective.
Indeed, if $R_v=M_{r_v}(\Delta_v)$ with $\Delta_v$ a division algebra, the positive generator of $K_0(R)$ maps to $r_v$ times the positive generator of $K_0(R_v)$.
Thus
\begin{equation}\label{eq:paper-lca-positive-row}
 X_R\xrightarrow{\lambda_R}B_R\longrightarrow Y_R
\end{equation}
is also a fiber sequence.
The map $K_1(R)\to K_1(\operatorname{AB}_R)$ is injective as well: for $R=D$, this follows from $\SK_1(D)=0$~\cite{Wang1950} and $F^\times\hookrightarrow F_v^\times$; for $R=F$, it follows from the injectivity of the diagonal map.
Hence applying $\tau_{\ge2}$ to all three terms of~\eqref{eq:paper-lca-positive-row} again gives a fiber sequence \cite[Construction 8.2 and Remark 8.3]{BHRAdeles}.

\subsection{Local--global comparison from a common integral model}
\label{subsec:paper-lca-comparison}
The two localization sequences are compared by means of a common integral model of $D$.
The difficulty is that the global comparison and the local comparisons are constructed separately, so the square relating them must be produced at the level of spectra and not only on homotopy groups.

\begin{lemma}
\label{lem:paper-lca-common-model}
Let $S_0$ be a finite set of finite places containing the finite places at which $D$ is ramified and the places above primes dividing $\deg D$, large enough that $D$ extends to an Azumaya algebra $\mathcal A$ over $U_0=\Spec\os_{F,S_0}$.
The geometric rank of a class over $U_0$ means its geometric rank at the generic point.
Let $\ell$ be a prime dividing $\deg D$, with $\ell\ne2$ if there is real ramification.
At a ramified real place $v$, where $2$ is inverted, the class $[\mathbb H]/2$ serves as a rank-$1$ class; here $\mathbb H$ is the division algebra in the Brauer class of $D_v$.
\begin{enumerate}
\item One can choose
\[
 x_{\ell,U_0}\in\pi_0L_{K(1),\ell}K(\mathcal A)
\]
of geometric rank $1$, with the prescribed rank-$1$ Morita class at every split real place.
\item For $x=x_{\ell,U_0}$, the multiplication map
\[
 m_x\colon L_{K(1),\ell}K(U_0)\longrightarrow L_{K(1),\ell}K(\mathcal A)
\]
is an equivalence.
\end{enumerate}
\end{lemma}
\begin{proof}
Apply the lifting of Proposition~\ref{prop:global-rank} over $U_0$, where $\ell$ is invertible.
Restriction of cohomology with finite coefficients to the real places is an isomorphism in degrees at least three \cite[I, Theorem 4.10(c)]{MilneADT}.
For degree two, choose a place $v\mid\ell$.
By local duality $H^2(F_v,M)\to H^0(F,M^D)^\vee$ is surjective, so by Poitou--Tate the obstruction for any tuple at the real places can be cancelled at $v$, which gives surjectivity onto the real components in degree two \cite[I, Corollary 4.16]{MilneADT}.
The long exact sequence of the real-relative fiber therefore shows that the real-relative sections functor has cohomological dimension at most two.
The only geometric homotopy sheaves with $\Z/\ell^a$ coefficients are $\Z/\ell^a(b)$ in even degrees $2b$.
If no values at real places need to be prescribed, the obstruction from the rank edge map lies in $H^3(U_0,\Z/\ell^a(1))=0$.
In the real-relative case, compute the obstruction to the prescribed tuple of real rank-$1$ classes using relative sections.
The only term remaining in total degree $-1$ is the weight-zero $H^1_{\mathrm{rel}}$, and we write $\delta$ for the second arrow in
\[
 H^0(U_0,\Z/\ell^a)\longrightarrow
 \prod_{v\ \mathrm{real}}H^0(F_v,\Z/\ell^a)
 \longrightarrow H^1_{\mathrm{rel}}(U_0,\Z/\ell^a).
\]
The obstruction $\delta((1)_v)$ vanishes because $(1)_v$ is the restriction of the diagonal element $1$.
Thus for each $a$ there is an element with the required rank and prescribed real values.

Let $E_a$ be the periodic \'etale hypersheaf $V\mapsto(L_{K(1),\ell}K(\mathcal A|_V))/\ell^a$; since $\ell$ is invertible on $U_0$, its global sections are $(L_{K(1),\ell}K(\mathcal A))/\ell^a$.
The rank is a section of a constant sheaf on the connected scheme $U_0$, so rank one at the generic point means rank one at every geometric point.
The fiber of the map from $\Gamma(U_0,\tau_{\ge0}E_a)$ to the pullback prescribing the global rank and real components of the same rank is $R\Gamma_{\mathrm{rel}}(U_0,\tau_{\ge2}E_a)$; by relative cohomological dimension at most two, its $\pi_{-1}$ vanishes and its $\pi_0$ is $H^2_{\mathrm{rel}}(U_0,\Z/\ell^a(1))$.
Thus any discrepancy with the previous choice under coefficient reduction $\ell^{a+1}\to\ell^a$ can be corrected by the weight-one $H^2$ (the relative $H^2$ in the real-relative case).
The next $H^3$ in the coefficient sequence vanishes, so reduction on this finite $H^2$ is surjective, and we can choose compatible lifts inductively.
Finiteness of the finite-coefficient groups over a finite $S_0$ gives the Mittag--Leffler condition, so the Milnor $\varprojlim^1$ vanishes, and $x_{\ell,U_0}$ is the inverse limit of this compatible system.

For $v\notin S_0$, the algebra $\mathcal A_v$ splits and $v\nmid\ell$.
Rigidity and the calculation over finite residue fields give
\[
 \pi_0L_{K(1),\ell}K(\os_v)=\Z_\ell.
\]
Hence $x_{\ell,U_0}|_v$ is precisely the rank-$1$ Morita class.
The same holds at infinite places: by prescription at split real places and by the rank calculation at complex places.
At a ramified real place, for odd $\ell$, geometric rank identifies $\pi_0$ of the height-one localization with $\Z_\ell$, and the rank-$1$ class in this group is $[\mathbb H]/2$.
We have $m_{[\mathbb H]/2}^{-1}=N_{\mathbb H/\mathbb R}/2$.
Thus, with $x=x_{\ell,U_0}$, the class $\beta_\ell=N_{\mathcal A/U_0}(x)-\deg D$ used below also restricts to zero at every infinite place.
Since $\mathcal A$ is \'etale locally a matrix algebra over $U_0$ and $\ell$ is invertible on $U_0$, the proof of Lemma~\ref{lem:paper-etale-rank-one} applies over $U_0$ and proves~(2).
The multiplication map $m_x$ and its inverse are constructed by restriction from the common model, so we obtain at the same time the homotopies expressing compatibility with restriction and the coherence of their composites.
\end{proof}

The single element $x_{\ell,U_0}$ and its base changes allow the completed comparisons over the global field, its completions at finite places, and the integer rings at $v\notin S_0$ to be chosen compatibly with restriction maps.
At $v\notin S_0$ and at split infinite places, these are the usual Morita comparisons.
At ramified real places, Morita invariance identifies the comparison with the normalized norm $N_{\mathbb H/\mathbb R}/2$; in the coordinates of the central simple algebra $D_v$ of degree $\deg D$ this is $N_{D_v/F_v}/\deg D$, since $N_{D_v/F_v}=(\deg D/2)N_{\mathbb H/\mathbb R}$.

We also fix Morita duals of these classes.
Let $\ell$ be a prime, assumed odd if $D$ is ramified at a real place, and put $L_\ell=L_{K(1),\ell}$.
If $\ell\mid\deg D$, let $x$ be the class $x_{\ell,U_0}$ of Lemma~\ref{lem:paper-lca-common-model}; if $\ell\nmid\deg D$, put $x=[\mathcal A]/\deg D$, which is defined $\ell$-locally and has geometric rank $1$.
For the Morita pairing
\[
 \mu\colon L_\ell K(\mathcal A)\otimes_{L_\ell K(U_0)}L_\ell K(\mathcal A^{\mathrm{op}})\longrightarrow L_\ell K(U_0),
\]
the map $\mu(x,-)$ is an equivalence, and we put $y=\mu(x,-)^{-1}(1)$, the \emph{Morita dual} of $x$.
If $\ell\mid\deg D$, then $\ell$ is invertible on $U_0$; \'etale locally on $U_0$ the tensor product of a simple right module and a simple left module is free of rank one, and $\mu(x,-)$ becomes multiplication by a unit, so the assertion follows by \'etale hyperdescent.
If $\ell\nmid\deg D$, we use the free-module argument instead.
The underlying vector bundle of $\mathcal A$ has trivial determinant, since \'etale locally $\mathcal A\simeq\operatorname{End}(E)$ and $\det\operatorname{End}(E)\simeq\os$ canonically, a trivialization that descends because $\mathrm{PGL}_n$ has no nontrivial characters.
With $d=\deg D$, let $i$ and $N$ denote extension of scalars to $\mathcal A$ and the norm.
Then $N\circ i$ is multiplication by $[\mathcal A]$ and $i\circ N$ by $[\mathcal A^\vee]$, the latter by the Morita equivalence $\mathcal A^e\simeq\operatorname{End}(\mathcal A)$; by the triviality of the determinant, both are $d^2$ in $K_0(U_0)$.
Hence $i/d$ and $N/d$ are mutually inverse, and $\mu([\mathcal A]/d,[\mathcal A^{\mathrm{op}}]/d)=[\mathcal A]/d^2=1$; on the opposite side $m_{[\mathcal A^{\mathrm{op}}]/d}$ is also an equivalence, so $\mu(x,-)$ is an equivalence and this element is the Morita dual, that is, $y=[\mathcal A^{\mathrm{op}}]/d$.
In both cases, $\mu(x,y)=1$, the map $m_y\colon L_\ell K(U_0)\to L_\ell K(\mathcal A^{\mathrm{op}})$ is an equivalence, and $\mu(-,y)\simeq m_x^{-1}$.
Over $F$, its completions, and the integer rings, we use the base changes of $x$ and $y$.

Let $g_D\colon X_D\to X_F$ be the comparison of Theorem~\ref{thm:paper-global-comparison} whose completed components are $m_x^{-1}$ for the base changes to $F$ of the classes $x_{\ell,U_0}$ ($\ell\mid\deg D$) and whose component with $\deg D$ inverted is $N_{D/F}/\deg D$.
The only remaining choice is the gluing homotopy in the arithmetic square, which we fix; it need not be compatible with the local comparisons, which is why the square below may have to be corrected.
Write $\nu_{j,D_v}$ for the isomorphisms of Theorem~\ref{thm:paper-local-norm} for $D_v$ at finite places $v$.
At infinite places, we use the same notation for the Morita comparison above, or for $N_{\mathbb H/\mathbb R}/2$ at a ramified real place.

\begin{theorem}
\label{thm:paper-lca-comparison}
One can choose a homotopy-commutative diagram
\begin{equation}\label{eq:paper-lca-comparison-square}
\begin{array}{ccccc}
 X_D&\xrightarrow{\lambda_D}&B_D&\longrightarrow&Y_D\\
 {\scriptstyle g_D}\downarrow&&
 {\scriptstyle h_D}\downarrow\,\simeq&&{\scriptstyle c^{\mathrm{lc}}_D}\downarrow\\
 X_F&\xrightarrow{\lambda_F}&B_F&\longrightarrow&Y_F
\end{array}
\end{equation}
whose rows are fiber sequences and whose middle vertical map is an equivalence.
Moreover, in every degree $j\ge1$, one can arrange that
\begin{equation}\label{eq:paper-lca-all-positive-norms}
 (h_D)_j=\prod'_v\nu_{j,D_v}.
\end{equation}
The dependence of $h_D$ on the choice of $g_D$ is described in Remark~\ref{rem:paper-lca-comparison-choices}.
\end{theorem}
The proof of Theorem~\ref{thm:paper-lca-comparison} is given after Lemma~\ref{lem:paper-lca-regulator-retract} below.

\begin{remark}\label{rem:paper-lca-comparison-choices}
If the global comparison is chosen simultaneously from the common model, $h_D$ itself can be taken to be the restricted product of the local spectrum comparisons.
If this gluing homotopy of $g_D$ is kept fixed, a correction preserving~\eqref{eq:paper-lca-all-positive-norms} may be needed; it is constructed in the proof below.
\end{remark}

We first construct, from the common integral model, a global comparison $g_c$ compatible with the local comparisons (Lemmas~\ref{lem:paper-local-d-arithmetic}--\ref{lem:paper-lca-compatible-gluing}), and then use a rational splitting (Lemma~\ref{lem:paper-lca-regulator-retract}) to return from $g_c$ to the given $g_D$.
Throughout, commutativity of squares of spectra is established by explicit homotopies and is not deduced from equality on homotopy groups.

For a positive integer $n$ and a spectrum $T$, put
\[
 T^\wedge_n=\prod_{\ell\mid n}T^\wedge_\ell.
\]
The arithmetic square for this finite set of primes is
\begin{equation}\label{eq:paper-finite-d-arithmetic}
 T\simeq T[1/n]\mathop{\times}_{T^\wedge_n[1/n]}T^\wedge_n,
\end{equation}
as one verifies by rationalization and comparison modulo each prime; the spectrum $T^\wedge_n[1/n]$ is rational.

\begin{lemma}
\label{lem:paper-local-d-arithmetic}
Let $k$ be a $p$-adic local field, let $D'$ be a central simple algebra over $k$ of degree $e$, not necessarily a division algebra, and let $n$ be a positive multiple of $e$.
One can choose the equivalence of Theorem~\ref{thm:paper-local-spectrum} so that it induces all the isomorphisms of Theorem~\ref{thm:paper-local-norm}, using a prescribed geometric rank-$1$ class at each $\ell\mid n$ and keeping the component with $n$ inverted equal to $N_{D'/k}/e$.
\end{lemma}
\begin{proof}
Put $X_{D'}=\tau_{\ge1}K(D')$ and $X_k=\tau_{\ge1}K(k)$.
If $D'\simeq M_r(\Delta)$, then under Morita equivalence $N_{D'/k}=r\,N_{\Delta/k}$, so Theorem~\ref{thm:paper-local-norm} gives $\pi_jN_{D'/k}=r\deg\Delta\,\nu_{j,\Delta}=e\,\nu_{j,D'}$ in all positive degrees; thus the factor is the degree $e$ of $D'$, not its index.

The prescribed completed comparisons and $N_{D'/k}$ agree on the overlap after dividing by $e$, by the projection formula and the vanishing of products with classes of geometric rank $0$ (Lemma~\ref{lem:paper-rank-zero-product}).
We choose a comparison homotopy in the category of rational spectra and use~\eqref{eq:paper-finite-d-arithmetic} to obtain an initial equivalence $g^0\colon X_{D'}\to X_k$.
The odd-degree calculation and the splitting-transfer and Bockstein argument in the proof of Theorem~\ref{thm:paper-local-spectrum} show that $g^0$ agrees with the prescribed $\nu_{j,D'}$ in odd degrees and on torsion in integral homotopy.
The components with $n$ inverted agree, both being $N_{D'/k}/e$, and in positive even degrees the difference on the completed components factors through a rational group into a finite group, hence vanishes.
Any remaining discrepancy is confined to positive even degrees $j$ and factors through the torsion-free quotient $K_j(D')/K_j(D')_{\mathrm{tors}}$, which is a rational vector space.

By exactness of the arithmetic fiber sequence, the target of this discrepancy is identified with
\[
 \begin{gathered}
 V_j/(\Gamma_{n,j}+\Lambda_j),\qquad V_j=\pi_{j+1}(X_k)^\wedge_n[1/n],\\
 \Gamma_{n,j}=\im\bigl(\pi_{j+1}X_k[1/n]\to V_j\bigr),\\
 \Lambda_j=\im\bigl(\pi_{j+1}(X_k)^\wedge_n\to V_j\bigr).
 \end{gathered}
\]
Since $j+1\ge3$ is odd, the completed groups at $\ell\ne p$ are finite, but $\pi_{j+1}(X_k)^\wedge_p$ is a finitely generated $\Z_p$-module.
Thus $V_j=0$ if $p\nmid n$, but if $p\mid n$, then
\[
 V_j=\pi_{j+1}(X_k)^\wedge_p\otimes_{\Z_p}\Q_p,\qquad
 \Lambda_j=\im(\pi_{j+1}(X_k)^\wedge_p\to V_j)=\pi_{j+1}(X_k)^\wedge_p/(\pi_{j+1}(X_k)^\wedge_p)_{\mathrm{tors}}
\]
and $\Lambda_j$ is a full $\Z_p$-lattice in $V_j$.
Since $\Gamma_{n,j}$ is $p$-divisible, Lemma~\ref{lem:paper-rational-correction}(3) gives
\[
 \operatorname{Ext}^1_{\Z}(\Q,\Gamma_{n,j}+\Lambda_j)=0.
\]
Lemma~\ref{lem:paper-rational-correction}(1), applied to~\eqref{eq:paper-finite-d-arithmetic}, now corrects $g^0$ to an equivalence $g$ that induces the prescribed isomorphisms, and the correction preserves the prescribed components.
\end{proof}

\begin{lemma}
\label{lem:paper-arithmetic-arrow-gluing}
Suppose given maps $a_R\colon X_R\to B'_R$ ($R=D,F$), an equivalence $h'\colon B'_D\to B'_F$, maps $u_{\deg D}\colon X_D[1/\deg D]\to X_F[1/\deg D]$ and $u_\ell\colon(X_D)^\wedge_\ell\to (X_F)^\wedge_\ell$ for each prime $\ell\mid\deg D$, and homotopies of squares compatible with $h'$.
We take the difference between the two comparisons on the overlap, including these homotopies of squares, and obtain a class
\[
 \Xi\in
 \bigl[X_D[1/\deg D],\fib
       \{(X_F)^\wedge_{\deg D}[1/\deg D]\to (B'_F)^\wedge_{\deg D}[1/\deg D]\}\bigr].
\]
If $\Xi=0$, one can choose $g_c\colon X_D\to X_F$ with the prescribed arithmetic components and homotopies of squares.
\end{lemma}
\begin{proof}
Apply~\eqref{eq:paper-finite-d-arithmetic} objectwise in the arrow category $\mathrm{Ar}=\operatorname{Fun}(\Delta^1,\mathrm{Sp})$.
Limits in $\mathrm{Ar}$ are computed objectwise, and the mapping spaces of $\mathrm{Ar}$ are given by the homotopy pullback
\[
 \operatorname{Map}_{\mathrm{Ar}}(a_D,a_F)
 \simeq\operatorname{Map}(X_D,X_F)
 \mathop{\times}_{\operatorname{Map}(X_D,B'_F)}
                         \operatorname{Map}(B'_D,B'_F).
\]
We fix the component $B'_D\to B'_F$ to be $h'$; the resulting homotopy fiber, if nonempty, is a torsor for $\operatorname{Map}(X_D,\fib(a_F))$.
By taking the difference between the two comparisons on the overlap, we pass from this torsor without a chosen base point to the zero fiber and obtain the class $\Xi$ above.
A nullhomotopy of $\Xi$ therefore supplies the homotopy needed to glue the two arithmetic components; we choose one and take the pullback, which preserves the prescribed component homotopies.
\end{proof}

We fix the following choices, which are used in Lemmas~\ref{lem:paper-lca-relative-cohomology}--\ref{lem:paper-lca-compatible-gluing}.
For each $v\in S_0$, apply Lemma~\ref{lem:paper-local-d-arithmetic} to $D'=D_v$ with $n=\deg D'=\deg D$; thus the prescribed primes are all $\ell\mid\deg D$, even when the local index is smaller.
Fix the completed components to be those determined by $x_{\ell,U_0}|_v$, and take $N_{D_v/F_v}/\deg D$ on the component with $\deg D$ inverted.
If there is real ramification, the completed components are taken at odd primes $\ell$ only.
We use the integral Morita comparisons at the remaining places and $N_{D_v/F_v}/\deg D$ at ramified real places, which is $N_{\mathbb H/\mathbb R}/2$ under Morita invariance.
The comparison homotopies from integer rings to fields are obtained by base change from the choices over the integer rings.
The products at each finite stage and their transition maps therefore give an equivalence $h_0\colon B_D\to B_F$.

We write $B'_R$ for the factor of $B_R$ obtained by omitting the finitely many field components in $S_0$,
\[
 B'_R=\tau_{\ge1}\varinjlim_{S\supseteq S_0}\left(
 \prod_{v\mid\infty}K(R_v)\times\prod_{v\in S\setminus S_0}K(R_v)
 \times\prod_{v\notin S,\ v\nmid\infty}K(\mathfrak R_v)\right).
\]
We write $a_R\colon X_R\to B'_R$ for the diagonal map and $h'\colon B'_D\to B'_F$ for the comparison between these factors.
With $h'$ fixed, we glue the global comparison using $u_{\deg D}=N_{D/F}/\deg D$ and the $u_\ell$ obtained from $m_x^{-1}$ on the common model.
On the integer-ring factors, the field factors at finite places $v\notin S_0$, and the split infinite-place factors of $B'_D$, the norm for the degree-$\deg D$ matrix algebra is $\deg D$ times the Morita comparison.
At a ramified real place it is $(\deg D/2)N_{\mathbb H/\mathbb R}$.
Thus the equality
\[
 h'[1/\deg D]\,a_D[1/\deg D]\simeq a_F[1/\deg D]\,(N_{D/F}/\deg D)
\]
of maps of spectra follows from natural isomorphisms of these exact functors.
In what follows, exclude $\ell=2$ when there is real ramification.
We take the homotopies of the squares relating $u_{\deg D}$ and $u_\ell$ to $h'$ from the projection formula for the tensor product and from base change of the common model.

For each $\ell\mid \deg D$, the element
\[
 \beta_\ell=N_{\mathcal A/U_0}(x_{\ell,U_0})-\deg D\in\pi_0L_{K(1),\ell}K(U_0)
\]
vanishes on every local component of $B'_F$.
We take these nullhomotopies from the two specified square homotopies: on each integer-ring component, their difference, multiplied by $\deg D$, precomposed with $m_x$ and evaluated at the unit, is a nullhomotopy of $\beta_\ell$, and on the field components at finite places we use its base change; at the infinite places we restrict the same two squares directly along $U_0\to F_v$.
The comparison multiplied by $\deg D$ lies on the $\ell$-completed side because of the natural norm identifications $(\deg D)u_{\deg D}=N$ and $(\deg D)h'_v\simeq N_v$; we subtract the rank-$1$ side from the norm side.
This gives nullhomotopies compatible with changes of the finite set $S\supset S_0$.
Thus $\beta_\ell$ lifts to an element $\beta_{\ell,\mathrm{rel},S}$ of
\[
 \mathcal F_\ell^S=\fib\!\left(
 L_{K(1),\ell}K(U_S)\longrightarrow \Pi_\ell^S\right),
\]
where
\[
 \Pi_\ell^S=
 \prod_{v\mid\infty}L_{K(1),\ell}K(F_v)
 \times\prod_{v\in S\setminus S_0}L_{K(1),\ell}K(F_v)
 \times\prod_{v\notin S,\ v\nmid\infty}
                         L_{K(1),\ell}K(\os_v).
\]
The relative weight-zero $H^0$ vanishes by injectivity of the diagonal restriction to the infinite places.
The odd geometric layers also vanish, so this lift has descent filtration at least $2$.

\begin{lemma}\label{lem:paper-lca-relative-cohomology}
Let $S\supseteq S_0$ be a finite set of finite places, where all places above $\ell$ lie in $S_0$, and assume $\ell\ne2$ if there is real ramification.
Define the relative continuous cohomology $H^i_{\mathrm{rel},S}(\Z_\ell(t))$ by the fiber of the map from the global complex over $U_S$ to the product of the local complexes in $\Pi_\ell^S$, in which each infinite place occurs exactly once.
Then
\[
 H^i_{\mathrm{rel},S}(\Z_\ell(t))=0\quad(i>3),\qquad
 H^3_{\mathrm{rel},S}(\Z_\ell(t))\text{ is finite}\quad(t\ge2).
\]
\end{lemma}
\begin{proof}
On the global side for finite $S$, inverse limits can be taken using finiteness of finite-coefficient groups; on the local side, derived inverse limits commute with products.
We use Milne's global and local duality.
The $H^i$ of the integer-ring factors vanish for $i\ge2$, and for $v\in S\setminus S_0$ one has $v\nmid\ell$, so
\[
 H^2(F_v,\Z_\ell(t))=
 \Z_\ell/(q_v^{t-1}-1)\Z_\ell\quad(t\ge2)
\]
is finite.
Restriction from global cohomology to the real places is an isomorphism in degrees at least three \cite[I, Theorem 4.10(c)]{MilneADT}.
The relative long exact sequence therefore expresses relative $H^3$ as the cokernel of a map into the direct sum of the finitely many finite local groups $H^2$ above and the finite groups at the real places, and higher relative degrees vanish \cite[I, Theorem 4.10 and Corollary 4.16]{MilneADT}.
\end{proof}

Here and below, the \emph{descent filtration} is the filtration induced by the relative geometric Postnikov tower; its $s$th graded piece in degree $j$ is a subquotient (the $E_\infty$-term) of $H^s_{\mathrm{rel},S}(\Z_\ell(t))$ with $2t-s=j$.
By Lemma~\ref{lem:paper-lca-relative-cohomology}, this tower converges in each fixed degree, and filtrations four and higher contribute nothing.

\begin{lemma}\label{lem:paper-lca-obstruction-vanishing}
For the common-model classes $x_{\ell,U_0}$ for all primes $\ell\mid\deg D$, with $\ell\ne2$ if there is real ramification, the comparison $h'$, the relative lifts $\beta_{\ell,\mathrm{rel},S}$ and the homotopies fixed above, the obstruction class $\Xi$ of Lemma~\ref{lem:paper-arithmetic-arrow-gluing} vanishes.
\end{lemma}
\begin{proof}
To compute $\Xi$ on the common model, fix $\ell$ and put $x=x_{\ell,U_0}$.
On the overlap, with $d=\deg D$, we have $d(u_d-u_\ell)\simeq(m_{N_{D/F}(x)}-d)\,m_x^{-1}=m_{\beta_\ell}m_x^{-1}$ by the projection formula, and the nullhomotopies of $\beta_\ell$ fixed above make this a map into the relative fiber.
Over finite $U_S$, the projection formula $N_{\mathcal A|_{U_S}/U_S}m_x\simeq m_{N(x)}$ leads us to consider the relative map on the common model
\begin{equation}\label{eq:paper-lca-relative-defect}
 m_{\beta_{\ell,\mathrm{rel},S}}m_x^{-1}.
\end{equation}
We use this expression as a map from $\tau_{\ge1}K(\mathcal A|_{U_S})$, so it factors through $\tau_{\ge1}\mathcal F_\ell^S$; by the choice of the nullhomotopies of $\beta_\ell$, it represents $d\,\Xi_\ell$ as a map into the relative fiber, not merely as a formula, where $\Xi_\ell$ is the $\ell$-component of $\Xi$, and multiplying by $d^{-1}$ on the overlap gives $\Xi_\ell$.
Since finite Moore coefficients commute with products and filtered colimits, we compare the common-model diagrams with finite coefficients in positive degrees for each finite $S$, take the colimit in the completed category, and then pass to the overlap.
Each global class used to test the obstruction is computed by returning to the relative spectrum $\mathcal F_\ell^S$ at a finite $S$ before this passage.
Since the target is rational, the source of $\Xi$ can be replaced by $X_{D,\Q}$.

A global class of positive odd degree $j=2m-1$ comes from some $U_S$.
Since $j$ is odd, the geometric layer in degree $j$ vanishes, so the completed image of this class has descent filtration at least $1$.
The map $m_x^{-1}$ is a map of \'etale modules and preserves the descent filtration.
The relative fiber $\mathcal F_\ell^S$ is a module over the global side.
For finite $S$, the class $\beta_{\ell,\mathrm{rel},S}$ lies in relative geometric Postnikov stage at least $2$, and after $m_x^{-1}$ the class of odd degree $j$ lifts to absolute geometric stage at least $j+1$.
Their product in~\eqref{eq:paper-lca-relative-defect} therefore lies in relative stage at least $j+3$, that is, in descent filtration at least $3$, and the only possible term remaining in degree $j=2m-1$ is $H^3_{\mathrm{rel},S}(\Z_\ell(m+1))$.
This group is finite because $m+1\ge2$, so it vanishes after rationalization.
In positive even degrees the source is $K_{2m}(D)\otimes\Q=0$, since $N_{D/F}i_D=i_DN_{D/F}=(\deg D)^2$ gives $K_{2m}(D)\otimes\Q\cong K_{2m}(F)\otimes\Q$, which vanishes by \cite{Borel}.
Integral classes generate the rationalized homotopy groups, so $\Xi$ is zero on all homotopy groups, and hence $\Xi=0$ because maps between rational spectra are determined by their effect on homotopy groups.
\end{proof}

\begin{lemma}
\label{lem:paper-lca-compatible-gluing}
For the choices in Lemma~\ref{lem:paper-lca-common-model}, one can choose a comparison $g_c\colon X_D\to X_F$ and an equivalence $h_0\colon B_D\to B_F$ given by the restricted product of the local comparisons such that
\[
 h_0\lambda_D\simeq\lambda_Fg_c.
\]
The maps $g_c$ and $g_D$ have the same completed components and the same component with $\deg D$ inverted.
\end{lemma}
\begin{proof}[Proof of Lemma~\ref{lem:paper-lca-compatible-gluing}]
The equivalence $h_0$ was constructed above.
By Lemma~\ref{lem:paper-lca-obstruction-vanishing}, Lemma~\ref{lem:paper-arithmetic-arrow-gluing} gives $g_c$, and the gluing preserves $h'$ and the prescribed homotopies of completed squares.

Finally, we restore the finitely many places $v\in S_0$.
For $X_{F_v}=\tau_{\ge1}K(F_v)$, the two localization maps agree on the completed components and the component with $\deg D$ inverted.
The obstruction to lifting their difference, together with the homotopies on those prescribed components, lies in
\[
 [X_D,\Omega (X_{F_v})^\wedge_{\deg D}[1/\deg D]]=[X_{D,\Q},\Omega (X_{F_v})^\wedge_{\deg D}[1/\deg D]].
\]
The rationalized source has homotopy only in positive odd degrees, and $(X_{F_v})^\wedge_{\deg D}[1/\deg D]$ has zero homotopy in positive even degrees.
Therefore the group above vanishes; equivalently,
\[
 \pi_1\operatorname{Map}\bigl(X_D,(X_{F_v})^\wedge_{\deg D}[1/\deg D]\bigr)=0.
\]
Hence the difference on the overlap between the homotopies prescribed on the two arithmetic components can be removed, and the places in $S_0$ can be restored while retaining the prescribed projection-formula homotopies of the common model.
\end{proof}

\begin{lemma}\label{lem:paper-lca-regulator-retract}
The map $\lambda_{F,\Q}\colon X_{F,\Q}\to B_{F,\Q}$ is a split monomorphism.
\end{lemma}
\begin{proof}
In degree one, $F^\times\to\mathbb A_F^\times$ is injective and remains so after rationalization.
In positive even degrees, $K_{2m}(F)_\Q=0$.
Let $r_1$ and $r_2$ be the numbers of real and complex places of $F$, respectively.
For $m\ge2$, put $r_m=r_1+r_2$ if $m$ is odd, and $r_m=r_2$ if $m$ is even.
In odd degrees $2m-1\ge3$, the Borel regulator through the infinite-place terms,
\[
 K_{2m-1}(F)\longrightarrow
 \prod_{v\mid\infty}K_{2m-1}(F_v)\longrightarrow\mathbb R^{r_m},
\]
is injective on the rational global $K$-group \cite{Borel}.
Hence $\lambda_{F,\Q}$ is injective on every homotopy group.
Since the category of rational spectra is equivalent to the category of complexes over $\Q$ and short exact sequences of $\Q$-vector spaces split, we can choose
\[
 \sigma_F\colon B_{F,\Q}\longrightarrow X_{F,\Q},\qquad
 \sigma_F\lambda_{F,\Q}\simeq\mathrm{id}.
\]
\end{proof}

\begin{proof}[Proof of Theorem~\ref{thm:paper-lca-comparison}]
We take $g_c$ and $h_0$ from Lemma~\ref{lem:paper-lca-compatible-gluing}.
Since $g_c-g_D$ vanishes on each component of the arithmetic square, we choose $b_0\colon X_D\to\Omega (X_F)^\wedge_{\deg D}[1/\deg D]$ with $g_c-g_D=\partial_F b_0$, where $\partial_F$ is the boundary in the arithmetic fiber sequence.
The target is rational, so $b_0$ factors through $X_{D,\Q}$.
Let $f$ be the composite of this factorization with $\partial_F$.
Then
\[
 g_c-g_D=fq_D,\qquad
 f\colon X_{D,\Q}\longrightarrow X_F,\qquad q_Ff\simeq0,
\]
where $q_D$ and $q_F$ are the rationalization maps.
Here $q_Ff\simeq0$ because the composite $\Omega (X_F)^\wedge_{\deg D}[1/\deg D]\xrightarrow{\partial_F}X_F\to X_F[1/\deg D]$ vanishes and $q_F$ factors through $X_F\to X_F[1/\deg D]$.
Since $g_D[1/\deg D]=N_{D/F}/\deg D$ and $N_{D/F}i_D=i_DN_{D/F}=(\deg D)^2$, the map $g_{D,\Q}$ is an equivalence, and putting
\[
 t=f(g_{D,\Q})^{-1}\colon X_{F,\Q}\longrightarrow X_F,
 \qquad \varepsilon=tq_F
\]
gives $q_Ft=0$, $\varepsilon^2=0$, and $g_c=(1+\varepsilon)g_D$.

Using $\sigma_F$ from Lemma~\ref{lem:paper-lca-regulator-retract}, we define
\[
 \zeta=\lambda_Ft\,\sigma_Fq_{B_F},\qquad b=1+\zeta.
\]
The identities $\sigma_F\lambda_{F,\Q}=1$ and $q_Ft=0$ give
\[
 \zeta^2=0,\qquad
 b^{-1}=1-\zeta,\qquad
 b\lambda_F=\lambda_F(1+\varepsilon).
\]
Here we choose $b_0$ together with the prescribed nullhomotopies of $g_c-g_D$ on the arithmetic components, that is, in the fiber of the corresponding mapping spaces; these data give $b\lambda_F\simeq\lambda_F(1+\varepsilon)$ and hence the square for $h_D$.
The correction pair $(\varepsilon,\zeta)$ is the map of arrow categories induced by $(t,\lambda_Ft\sigma_F)$ and $\sigma_F\lambda_{F,\Q}\simeq1$, so it factors through $\lambda_{F,\Q}$; in particular, for every prime $p$, the height-one localization $L_{K(1),p}$ kills the correction together with the homotopy of the square.
Consequently, $h_D=b^{-1}h_0$ is an equivalence, and
\[
 h_D\lambda_D\simeq b^{-1}\lambda_Fg_c
 \simeq b^{-1}\lambda_F(1+\varepsilon)g_D
 \simeq\lambda_Fg_D.
\]
We use this homotopy to define the map $c^{\mathrm{lc}}_D$ between cofibers.

Finally, we check that this correction does not change the norms on homotopy groups.
For $j\ge1$, the source of $t$ is rational, so the image of $t_*$ in $K_j(F)$, or in $K_j(F)[1/2]$ under real ramification, is divisible and hence zero by Lemma~\ref{lem:paper-global-no-divisible}.
Thus $\zeta_*=0$ and $b_*=1$, so $h_D$ satisfies the same formula~\eqref{eq:paper-lca-all-positive-norms} as $h_0$.
Moreover, $\zeta$ factors through a rational spectrum and hence vanishes with all finite coefficients, so the completed components are unchanged as well.
\end{proof}

\begin{remark}\label{rem:paper-lca-choice}
Theorem~\ref{thm:paper-lca-comparison} is an existence result, and its construction involves choices.
We do not deduce that the correction maps $t$ or $\zeta$ vanish as maps of spectra merely because they vanish on homotopy groups.
On the other hand, since $g_c=(1+\varepsilon)g_D$, returning from $g_c$ to $g_D$ changes the comparison by a self-equivalence of the target and therefore preserves the comparison fiber and the splitting in degrees at least two.
When expressing $m_x^{-1}$ on a periodic completed component by the Morita pairing $\mu$, we use the Morita dual $y$ of $x$ fixed after Lemma~\ref{lem:paper-lca-common-model} and its base changes.
This normalizes the comparison map more strongly than merely prescribing geometric rank.
The homotopies on the completed components of~\eqref{eq:paper-lca-comparison-square} are also chosen to be the projection-formula homotopies associated with base change of the common model and evaluation against $y$.
Corrections factoring through rational spectra vanish under height-one localization, so the height-one localization of~\eqref{eq:paper-lca-comparison-square} is preserved after correction.
The map $c^{\mathrm{lc}}_D$ is the cofiber map of the specified square.
Independence of the choice of commutativity homotopy is not asserted.
\end{remark}

\begin{lemma}
\label{lem:paper-lca-periodic-evaluation}
Let $p$ be a prime, assumed odd if there is real ramification.
Let $x$ and $y$ be the class at $p$ on the common model and its Morita dual, fixed after Lemma~\ref{lem:paper-lca-common-model}.
Put $L=L_{K(1),p}$, and let $h_0$ be the equivalence of Lemma~\ref{lem:paper-lca-compatible-gluing}.
Under the equivalences $LB_R\simeq L\Theta_{R,p}$ of the proof, $Lh_0$ is the $L$-localization of the evaluation map $\Theta_{D,p}\to\Theta_{F,p}$ given componentwise by $\mu(-,y|_v)\simeq m_{x|_v}^{-1}$.
Consequently $L$ applied to~\eqref{eq:paper-lca-comparison-square} is the map of localized rows given by evaluation against $y$, and $Lc^{\mathrm{lc}}_D$ is the induced map on cofibers.
\end{lemma}
\begin{proof}
First observe that $LT=0$ for every bounded-above spectrum $T$.
We choose a finite type $1$ spectrum with a positive-degree $v_1$ self-map \cite{HopkinsSmith1998}.
Smashing its telescope with $T$ gives a contractible spectrum, since in each fixed degree the terms of the resulting sequential colimit eventually vanish.
The smash product of $K(1)$ with this finite spectrum is a finite direct sum of suspensions of nonzero $K(1)$, and the self-map acts on it by an equivalence.
Hence $K(1)\otimes T=0$ and therefore $LT=0$.
This applies in particular to the part removed by taking a connective cover.

Let $\mathfrak R$ be a maximal order in the common integral model, and define, in spectra,
\[
 \Theta_{R,p}=\varinjlim_S\left(
 \prod_{v\mid\infty}LK(R_v)\times\prod_{v\in S}LK(R_v)
 \times\prod_{v\notin S,\ v\nmid\infty}LK(\mathfrak R_v)\right).
\]
We do not assume that this object is already local.
The natural comparison $B_R\to \Theta_{R,p}$ induces isomorphisms on mod-$p$ homotopy groups in degrees $\ge2$.
Indeed, the finite-coefficient comparisons over nonarchimedean local fields, integer rings, $\mathbb R$, and $\mathbb C$ all hold in this uniform range.
For the nonarchimedean field components, use Lemma~\ref{lem:paper-local-compare}.
Restricting the colimit to $S\supset S_0$ makes every algebra in an integer-ring component split.
At $v\nmid p$, use Henselian rigidity and the finite-residue-field calculation.
At $v\mid p$, the localization sequence
\[
 K(k_v)\longrightarrow K(\os_v)\longrightarrow K(F_v)
\]
together with $K(k_v)/p\simeq H\F_p$ and $LK(k_v)=0$, shows that the integer-ring and field comparisons agree in degrees at least two.
For complex components, use Suslin's finite-coefficient comparison.
For real components at $p=2$, use the comparison of real algebraic $K$-theory with $KO$; at odd primes, restriction and transfer for $\mathbb R\subset\mathbb C$, together with invertibility of $2$, reduce to the complex comparison.
At ramified real components, $p$ is odd and $N_{\mathbb H/\mathbb R}/2$ reduces the problem to the same calculation over $\mathbb R$ \cite{QuillenFinite,ClausenMathewMorrow,SuslinClosed,KaroubiWeibelReal}.
Finite Moore spectra are dualizable and hence commute with products; the same range is preserved in the comparisons involving filtered colimits and positive covers.
The cofiber modulo $p$ is therefore bounded above.
The vanishing proved above and conservativity of reduction modulo $p$ on $K(1)$-local spectra give
\[
 LB_R\xrightarrow{\sim}L\Theta_{R,p}.
\]
For the conservativity assertion, if $T/p=0$, then $p$ is invertible on $T$ but acts by zero on $K(1)$, so $K(1)\otimes T=0$; if $T$ is also $K(1)$-local, then $T=0$.

Tensoring with perfect modules over the common model is compatible with all local factors and transition maps, so $\Theta_{R,p}$ can be constructed in the category of $LK(U_0)$-modules.
Evaluation against $y$ gives $\Theta_{D,p}\to\Theta_{F,p}$ with local components $m_{x|_v}^{-1}$.
Taking products and the colimit of the local homotopies used to construct $h_0$ gives a commutative square with the maps $B_R\to\Theta_{R,p}$.
Hence $Lh_0$ is the height-one localization of this evaluation.

Tensoring with a perfect left $\mathcal A$-module preserves the exact structures on discrete modules, adelic blocks, and locally compact modules, and commutes naturally with localization.
Hence tensoring with perfect modules over the common model also acts on the localization sequence of Theorem~\ref{thm:paper-lca-source}, giving a single diagram of the rows and their cofibers.
In the gluing above, we have prescribed the homotopies on the completed components of this diagram; if $p\nmid\deg D$, the relevant component is the one with $\deg D$ inverted, where $y=[\mathcal A^{\mathrm{op}}]/\deg D$ and we use the projection-formula homotopy for $N_{D/F}/\deg D$.
Lifting at the finitely many exceptional places preserves them, and the rational correction returning to the fixed $g_D$ vanishes under $L$.
Applying $L$ kills the overlap; for $p\mid\deg D$ the $p$-completed component survives and for $p\nmid\deg D$ the component with $\deg D$ inverted, so the whole square agrees with the prescribed evaluation.
Thus $Lc^{\mathrm{lc}}_D$ is evaluation against the same $y$.
\end{proof}

\subsection{The suspension of the comparison fiber and splitting in degrees at least two}
\label{subsec:paper-lca-fiber}
The global splitting in Section~\ref{sec:global-spectrum} passes to the cofibers of the comparison square.

\begin{theorem}
\label{thm:paper-lca-fiber-splitting}
Let $Z=\fib(g_D)$.
The comparison of Theorem~\ref{thm:paper-lca-comparison} satisfies
\begin{equation}\label{eq:paper-lca-fiber-suspension}
 \fib(c^{\mathrm{lc}}_D)\simeq\Sigma Z.
\end{equation}
There is also a noncanonical equivalence whose first projection is $\tau_{\ge2}c^{\mathrm{lc}}_D$,
\begin{equation}\label{eq:paper-lca-splitting}
 \tau_{\ge2}K(\operatorname{LCA}_D)
 \simeq\tau_{\ge2}K(\operatorname{LCA}_F)\oplus\Sigma Z.
\end{equation}
On the second summand, the localization boundary is the natural inclusion $\Sigma Z\to\Sigma\tau_{\ge2}K(D)$.
\end{theorem}
\begin{proof}
By the $3\times3$ lemma, the vertical fibers in~\eqref{eq:paper-lca-comparison-square} form a fiber sequence
\[
 Z\longrightarrow\fib(h_D)\longrightarrow\fib(c^{\mathrm{lc}}_D).
\]
Since $h_D$ is an equivalence, the middle term vanishes, giving \eqref{eq:paper-lca-fiber-suspension}.

We have checked after~\eqref{eq:paper-lca-positive-row} that applying $\tau_{\ge2}$ to all three terms still gives a fiber sequence.
By Theorem~\ref{thm:paper-global-split}, we may write
\[
 \tau_{\ge2}K(D)\simeq\tau_{\ge2}K(F)\oplus Z,
 \qquad\tau_{\ge2}g_D=\operatorname{pr}_1.
\]
After identifying the adelic terms by $\tau_{\ge2}h_D$, we check, including the homotopies, that $\lambda_D\simeq(\lambda_F,0)$.
Let $s$ be a section of $\tau_{\ge2}g_D$ and $i\colon Z\to\tau_{\ge2}K(D)$ the fiber inclusion, and use the decomposition given by $s\oplus i$.
Restricting the specified homotopy of the square along $s$ gives $\lambda_Ds\simeq\lambda_Fg_Ds\simeq\lambda_F$.
Restricting the same homotopy along $i$ and composing with the fiber nullhomotopy $g_Di\simeq0$ gives $\lambda_Di\simeq\lambda_Fg_Di\simeq0$.
Thus $\lambda_D\simeq(\lambda_F,0)$ in the arrow category, and this equivalence also identifies $c^{\mathrm{lc}}_D$ on cofibers with the first projection.
Taking cofibers gives
\[
 \cofib(\lambda_F,0)\simeq\cofib(\lambda_F)\oplus\Sigma Z,
\]
which proves~\eqref{eq:paper-lca-splitting}.
The final assertion follows from the definition of the cofiber boundary.
\end{proof}

\begin{corollary}
\label{cor:paper-lca-gem}
Let $e_\ell=\ell^{v_\ell(\deg D)}$ and $G_{\ell,j}=\pi_jZ_{(\ell)}$, as in Theorem~\ref{thm:paper-global-fiber-gem}.
Then
\[
 \fib(c^{\mathrm{lc}}_D)\simeq
 \bigoplus_{\ell\mid\deg D}\ \bigoplus_{j=4}^{e_\ell-3}
       \Sigma^{j+1}H(G_{\ell,j}),
\]
and, for $n\ge2$,
\[
 K_n(\operatorname{LCA}_D)\cong
 K_n(\operatorname{LCA}_F)\oplus
       \bigoplus_{\ell\mid\deg D}G_{\ell,n-1}.
\]
Here we set $G_{\ell,j}=0$ outside the stated range, and we omit $\ell=2$ from the sums if there is real ramification.
In particular, $c^{\mathrm{lc}}_D$ is an isomorphism in degrees $1$ through $4$, and the first possible difference is in degree $5$.
If $e_3=9$, then at $\ell=3$ the only possible differences are three finite groups, in degrees $5$, $6$, and $7$.
\end{corollary}
\begin{proof}
We combine Theorem~\ref{thm:paper-global-fiber-gem} with~\eqref{eq:paper-lca-fiber-suspension} and~\eqref{eq:paper-lca-splitting}.
The degree-one isomorphism follows from the long exact sequence of the fiber sequence $\fib(c^{\mathrm{lc}}_D)\to Y_D\to Y_F$.
No assertion is made that any of these finite groups is nonzero.
\end{proof}

\begin{theorem}\label{thm:paper-lca-k2}
For every central division algebra $D$ over $F$, there is an isomorphism
\[
 K_2(\operatorname{LCA}_D)\cong K_2(\operatorname{LCA}_F)
\]
with integral coefficients, also when $D$ is ramified at a real place.
In particular, \cite[Conjecture 10.2]{BHRAdeles} holds.
\end{theorem}
\begin{proof}
By \cite[Theorem 10.5]{BHRAdeles}, it suffices that, for $D$ and for every completion $D_v$, the reduced norm on $K_2$ is injective and its image is the kernel of the Hilbert symbols at the real places at which the algebra is ramified.
For $D$ this is Theorem~\ref{thm:paper-global-norm}, and for $D_v$ at a finite place $v$ it is Theorem~\ref{thm:paper-local-norm}.
At an infinite place at which $D$ splits, it holds by Morita equivalence.
At a real place at which $D$ is ramified, $D_v$ is Morita equivalent to the Hamilton quaternions $\mathbb H$, of index $2$.
By \cite[Corollary 2]{KahnLevine2010}, whose hypothesis is supplied by the norm residue isomorphism theorem \cite{Voevodsky}, $\Nrd_2$ is injective on $K_2(\mathbb H)$ with image the kernel of $[\mathbb H]\cup-\colon K_2(\mathbb R)\to H^4(\mathbb R,\Z/2)$.
Since $[\mathbb H]\cup(a)\cup(b)\ne0$ exactly when $a$ and $b$ are both negative, this kernel is the kernel of the Hilbert symbol.
\end{proof}

\subsection{Selmer \texorpdfstring{$K$}{K}-homology and id\`ele classes}\label{subsec:paper-lca-selmer}
We use Selmer $K$-homology to compare $c^{\mathrm{lc}}_D$ with Clausen's pro-$p$ Artin map and compute the map on id\`ele classes (Theorem~\ref{thm:paper-lca-artin}).

Let $p$ be a prime, assumed odd if $D$ is ramified at a real place.
Write $L=L_{K(1),p}$.
We write $\mathrm{lc}_\Z$ for the bounded derived category of second-countable locally compact abelian groups $\operatorname{LCA}_{\Z,\aleph}$, as in \cite[Definitions 0.2 and 3.2]{ClausenArtin}, $\omega=\omega_{K(1),p}$ for the $K(1)$-local dualizing object, $\mathbb D_p(T)=\operatorname{map}_{\mathrm{Sp}}(T,\omega)$ for the function spectrum.
For a small idempotent-complete $\Z$-linear stable $\infty$-category $\mathcal P$, Clausen's $p$-adic Selmer $K$-homology $dK^{\mathrm{Sel}}_p(\mathcal P)$ is the homotopy pushout of $\mathbb D_pK(\mathcal P)\leftarrow\mathbb D_p\mathrm{TC}(\mathcal P)\to d_{\mathrm{TC}}\mathrm{TC}(\mathcal P)$, where the left map is dual to the cyclotomic trace and $d_{\mathrm{TC}}$ is Clausen's duality for $\mathrm{TC}$ \cite[Definition 2.1]{ClausenArtin}; if $p$ is invertible in $\mathcal P$, it is equivalent to $\mathbb D_pK(\mathcal P)$ \cite[Proposition 2.11(1)]{ClausenArtin}.
We write $\mathbf j\in\pi_0dK^{\mathrm{Sel}}_p(\mathrm{lc}_\Z)$ for Clausen's fundamental class \cite[Definitions 1.4 and 4.2]{ClausenArtin}.
Under our convention of right $D$-modules, the perfect complexes used for evaluation are left $D$-modules, that is, objects of $\operatorname{Perf}(D^{\mathrm{op}})$.
Accordingly, define the Selmer objects by
\[
 \mathsf S_{D,p}=dK^{\mathrm{Sel}}_p(D^{\mathrm{op}}),
 \qquad
 \mathsf S_{F,p}=dK^{\mathrm{Sel}}_p(F).
\]
The evaluations below are performed on second-countable locally compact modules and transferred to the whole category using the $K$-equivalence of the inclusion in Lemma~\ref{lem:paper-lca-countability}.
Indeed, a second-countable right $D$-module $M_0$ defines an exact functor $\operatorname{Perf}(D^{\mathrm{op}})\to\mathrm{lc}_{\Z}$ by $P\mapsto M_0\otimes_D P$.
This correspondence induces an exact functor
\[
 D^b(\operatorname{LCA}_{D,\aleph})\longrightarrow
 \operatorname{Fun}_{\Z}
 \bigl(\operatorname{Perf}(D^{\mathrm{op}}),\mathrm{lc}_{\Z}\bigr).
\]
Pulling back Clausen's fundamental class along this evaluation gives
\[
 \operatorname{Art}_{D,p}\colon Y_D\longrightarrow\mathsf S_{D,p}.
\]
This is the composite of the Artin map for $\mathrm{lc}_{\operatorname{Perf}(D^{\mathrm{op}})}=\operatorname{Fun}_{\Z}(\operatorname{Perf}(D^{\mathrm{op}}),\mathrm{lc}_{\Z})$ with the $K$-theory of the exact functor above.
No equivalence between the stable categories themselves is required \cite[Definitions 3.2 and 4.2, Remarks 4.3--4.4]{ClausenArtin}.

\begin{theorem}\label{thm:paper-lca-artin}
Let $x$ be the class at $p$ on the common model and $y$ its Morita dual, fixed after Lemma~\ref{lem:paper-lca-common-model}.
Then $\rho_{D,p}=\mathbb D_p(m_y)\colon\mathsf S_{D,p}\to\mathsf S_{F,p}$ is an equivalence, and the comparison $c^{\mathrm{lc}}_D\colon Y_D\to Y_F$ of Theorem~\ref{thm:paper-lca-comparison} fits into a homotopy-commutative square
\begin{equation}\label{eq:paper-lca-artin-square}
 \begin{array}{ccc}
 Y_D&\xrightarrow{\operatorname{Art}_{D,p}}&\mathsf S_{D,p}\\
 {\scriptstyle c^{\mathrm{lc}}_D}\downarrow&&\downarrow{\scriptstyle\rho_{D,p}}\\
 Y_F&\xrightarrow{\operatorname{Art}_{F,p}}&\mathsf S_{F,p}.
 \end{array}
\end{equation}
Moreover, the restriction of $\operatorname{Art}_{D,p}$ to $\fib(c^{\mathrm{lc}}_D)\simeq\Sigma Z$ is zero.
In terms of Clausen's isomorphism $e_{F,p}\colon G_F^{\mathrm{ab},\wedge p} \xrightarrow{\sim}\pi_1\mathsf S_{F,p}$, the map on id\`ele classes is given, for $[(a_v)_v]\in\pi_1Y_D=\prod'_vK_1(D_v)/K_1(D)$, by
\begin{equation}\label{eq:paper-lca-reduced-artin}
 e_{F,p}^{-1}\,\pi_1(\rho_{D,p}\operatorname{Art}_{D,p})
 \bigl([(a_v)_v]\bigr)
 =
 \operatorname{rec}_{F,p}
 \bigl([(\Nrd_{D_v/F_v}(a_v))_v]\bigr).
\end{equation}
Here $a_v\in K_1(D_v)$ is assumed to come from the integral model at all but finitely many finite places.
The map $\operatorname{rec}_{F,p}$ is the global Artin map normalized so that a uniformizer at an unramified place maps to the geometric Frobenius, followed by the projection to the maximal pro-$p$ quotient, and $e_{F,p}$ is normalized as in \cite[Definition 2.37]{ClausenArtin}.
\end{theorem}

We prove Theorem~\ref{thm:paper-lca-artin} after two lemmas, which share the following choices.
Let $x\in\pi_0LK(D)$ and $y\in\pi_0LK(D^{\mathrm{op}})$ be the base changes to $F$ of the class at $p$ on the common model and its Morita dual, fixed after Lemma~\ref{lem:paper-lca-common-model}.
We also fix the homotopy making~\eqref{eq:paper-lca-comparison-square} commute, chosen as in Remark~\ref{rem:paper-lca-choice}, which is compatible with the localization sequences.

\begin{lemma}\label{lem:paper-lca-selmer-evaluation}
With these choices, $m_y\colon LK(F)\to LK(D^{\mathrm{op}})$ is an equivalence, and so is
\[
 \rho_{D,p}=\mathbb{D}_p(m_y)\colon\mathsf S_{D,p}\longrightarrow\mathsf S_{F,p}.
\]
The Selmer $K$-theory of the bilinear functor $(M_0,P)\mapsto M_0\otimes_DP$ defines a map $\Psi_y\colon LY_D\to LY_F$, evaluation against $y$, and $Lc^{\mathrm{lc}}_D\simeq\Psi_y$ as maps of localization sequences.
\end{lemma}
\begin{proof}
Put $M=LK(D)$ and $M'=LK(D^{\mathrm{op}})$.
Since the categories over a number field are $\Z[1/p]$-linear, \cite[Proposition 2.11]{ClausenArtin} gives
\[
 \mathsf S_{D,p}\simeq \mathbb{D}_p(M'),\qquad
 \mathsf S_{F,p}\simeq \mathbb{D}_p(LK(F)).
\]
By the choice of $y$, the map $m_y\colon LK(F)\to M'$ is an equivalence and $\mu(-,y)\simeq m_x^{-1}\colon M\to LK(F)$.
If $p\nmid \deg D$, then $y=[D^{\mathrm{op}}]/\deg D$ and $\mu(-,y)\simeq N_{D/F}/\deg D$, where $N_{D/F}$ is the norm used in Section~\ref{subsec:paper-lca-comparison}.
Thus $\rho_{D,p}=\mathbb D_p(m_y)$ is an equivalence.

In general, $y$ is not represented by an actual perfect module, so its action is defined on localized $K$-theory.
Put $\mathcal L_D=D^b(\operatorname{LCA}_{D,\aleph})$ and $\mathcal L_F=D^b(\operatorname{LCA}_{F,\aleph})$.
Tensoring with a finitely generated projective left $D$-module is defined on locally compact modules as a direct summand of a finite direct sum, and preserves second countability and exact sequences.
It therefore gives an exact bilinear functor
\[
 \mathcal L_D\otimes_F\operatorname{Perf}(D^{\mathrm{op}})
       \longrightarrow\mathcal L_F,
 \qquad (M_0,P)\longmapsto M_0\otimes_D P.
\]
The categories above and $\mathrm{lc}_{\operatorname{Perf}(D^{\mathrm{op}})}$ are $\Z[1/p]$-linear, so their $p$-completed $\operatorname{TC}$ vanishes, as for Selmer $K$-homology in \cite[Proposition 2.11]{ClausenArtin}.
For each of these categories $\mathcal P$, the definition of Selmer $K$-theory therefore gives
\[
 K^{\mathrm{Sel}}_p(\mathcal P)\simeq LK(\mathcal P).
\]
Thus the Selmer $K$-theory of this functor defines a pairing of spectra over $LK(F)$,
\[
 LK(\mathcal L_D)\otimes_{LK(F)}M'\longrightarrow LK(\mathcal L_F).
\]
Evaluation against $y$ gives $\Psi_y$ by substituting $m_y\colon LK(F)\to M'$ in the second variable.

Local maximal orders are hereditary and have no negative $K$-groups, so every factor in~\eqref{eq:paper-lca-restricted-spectrum}, and hence the filtered colimit, is connective.
Since $K(D)$ and $K(F)$ are connective, the localization sequences show that $K(\operatorname{LCA}_D)$ and $K(\operatorname{LCA}_F)$ are connective as well.
Taking positive covers thus removes only $H(K_0)$, which vanishes under height-one localization.
Using Lemma~\ref{lem:paper-lca-countability}, we therefore have $LK(\mathcal L_D)\simeq LY_D$ and $LK(\mathcal L_F)\simeq LY_F$.
With the homotopy making~\eqref{eq:paper-lca-comparison-square} commute fixed before this lemma, Lemma~\ref{lem:paper-lca-periodic-evaluation} gives, as maps of entire localization sequences,
\[
                         Lc^{\mathrm{lc}}_D\simeq \Psi_y.
\]
\end{proof}

\begin{lemma}\label{lem:paper-lca-artin-compatibility}
With the choices of Lemma~\ref{lem:paper-lca-selmer-evaluation},
\[
 \rho_{D,p}\operatorname{Art}_{D,p}^{L}
 \simeq\operatorname{Art}_{F,p}^{L}\Psi_y
 \simeq\operatorname{Art}_{F,p}^{L}Lc^{\mathrm{lc}}_D,
\]
where the superscript $L$ denotes factorization through height-one localization.
\end{lemma}
\begin{proof}
The Artin map of \cite[Definition 4.2]{ClausenArtin} is the composite of the Selmer $K$-theory evaluation pairing with the fundamental class $\mathbf j$.
We first define $\mathbf j_F\colon LK(D^b(\operatorname{LCA}_{F,\aleph}))\to\omega$ as the pullback of $\mathbf j$ along the forgetful functor to $\mathrm{lc}_\Z$, and we compare both sides after composing with the same $\mathbf j_F$.
Associativity of evaluation,
\[
 (M_0\otimes_D P)\otimes_F P'
 \simeq M_0\otimes_D(P\otimes_F P'),
\]
is a natural isomorphism of bilinear functors for this functor and the action of $\operatorname{Perf}(F)$, where $P'\in\operatorname{Perf}(F)$.
Passing to Selmer $K$-theory and pairing with $\mathbf j$, the two adjoints in question are
\[
 LY_D\otimes M'\longrightarrow LY_F\xrightarrow{\mathbf j_F}\omega
 \qquad\text{and}\qquad
 LY_F\otimes LK(F)\longrightarrow LY_F\xrightarrow{\mathbf j_F}\omega,
\]
where $M'=LK(D^{\mathrm{op}})$ and the first arrows are the evaluation pairing and the action of $LK(F)$.
Precomposed with $1\otimes m_y$ and with $\Psi_y\otimes1$, respectively, they agree as \emph{maps of spectra} $LY_D\otimes LK(F)\to\omega$, by balancedness of the evaluation pairing over $LK(F)$:
\[
 \operatorname{ev}\circ
   (\operatorname{Art}_{D,p}^{L}\otimes m_y)
 \simeq
 \operatorname{ev}\circ
   ((\operatorname{Art}_{F,p}^{L}\Psi_y)\otimes\mathrm{id}_{LK(F)}).
\]
By the internal-Hom adjunction and $\rho_{D,p}=\mathbb{D}_p(m_y)$,
\[
 \rho_{D,p}\operatorname{Art}_{D,p}^{L}
 \simeq\operatorname{Art}_{F,p}^{L}\Psi_y
 \simeq\operatorname{Art}_{F,p}^{L}Lc^{\mathrm{lc}}_D.
\]
The second equivalence is Lemma~\ref{lem:paper-lca-selmer-evaluation}.
\end{proof}

\begin{proof}[Proof of Theorem~\ref{thm:paper-lca-artin}]
Lemma~\ref{lem:paper-lca-artin-compatibility}, applied with the choices and the equivalence $\rho_{D,p}$ of Lemma~\ref{lem:paper-lca-selmer-evaluation}, gives $\rho_{D,p}\operatorname{Art}_{D,p}^{L}\simeq\operatorname{Art}_{F,p}^{L}Lc^{\mathrm{lc}}_D$.
Both Selmer objects are height-one local, so precomposing with $Y_D\to LY_D$ gives~\eqref{eq:paper-lca-artin-square}.
The choices of $x$, $y$, and the commutativity homotopy for the localization sequence are fixed simultaneously; independence of the choice of $\rho_{D,p}$ or $c^{\mathrm{lc}}_D$ is not asserted \cite[Definition 4.2 and Remarks 4.3--4.4]{ClausenArtin}.

We can check even more directly that the restriction of $\operatorname{Art}_{D,p}$ to $\fib(c^{\mathrm{lc}}_D)$ vanishes.
By Theorem~\ref{thm:paper-global-fiber-gem}, $Z^\wedge_p$ is a finite GEM with finite $p$-group coefficients.
Height-one localization is unchanged by $p$-completion and annihilates each Eilenberg--Mac Lane factor, so $L(\Sigma Z)=0$.
Consequently,
\[
 \operatorname{Map}_{\mathrm{Sp}}(\Sigma Z,\mathsf S_{D,p})
 \simeq
 \operatorname{Map}_{\mathrm{Sp}}(L\Sigma Z,\mathsf S_{D,p})
 \simeq *.
\]
The restriction of $\operatorname{Art}_{D,p}$ to $\fib(c^{\mathrm{lc}}_D)$ therefore vanishes, even independently of the choice of splitting.

Finally, since $\pi_0X_D=0$, the fiber sequence~\eqref{eq:paper-lca-positive-row} gives $\pi_1Y_D=\coker\bigl(K_1(D)\to\prod'_vK_1(D_v)\bigr)$, and similarly for $F$, so $\pi_1(c^{\mathrm{lc}}_D)$ is the map induced by $(h_D)_1$.
If there is real ramification, we first form the integral cokernel; since the target has $p$ odd, $2$ is invertible there and the formula extends uniquely to $\Z[1/2]$ coefficients.
Here $\nu_{1,D_v}=\Nrd_{D_v/F_v}$ for every $v$: at finite places by Theorem~\ref{thm:paper-local-norm}, at split infinite places because the Morita comparison on $K_1$ is the determinant, and at a ramified real place because $\pi_1N_{\mathbb H/\mathbb R}=2\Nrd_{\mathbb H/\mathbb R}$ on $K_1(\mathbb H)$.
Hence the case $j=1$ of~\eqref{eq:paper-lca-all-positive-norms} in Section~\ref{subsec:paper-lca-comparison} gives $\pi_1(c^{\mathrm{lc}}_D)([(a_v)_v])=[(\Nrd_{D_v/F_v}(a_v))_v]$.
On the field side, the class of an id\`ele in $\pi_1Y_F$ is represented by multiplication on $\mathbb A_F$, which is the description used in Clausen's comparison map \cite[Proposition 3.28(1) and Corollary 3.29(1)]{ClausenArtin}.
With $e_{F,p}$ defined as the negative of the Pontryagin-dual edge map, \cite[Theorem 0.9]{ClausenArtin} identifies $e_{F,p}^{-1}\circ\pi_1(\operatorname{Art}_{F,p})$ with $\operatorname{rec}_{F,p}$, as normalized in the theorem, on id\`ele classes \cite[Definition 2.37 and Lemma 4.5]{ClausenArtin}.
Combining this with~\eqref{eq:paper-lca-artin-square}, we obtain~\eqref{eq:paper-lca-reduced-artin}.
\end{proof}
